%========================================================================================
%  CAPITOLO 6 — STUDIO DELLE DEGENERAZIONI  (il cuore della tesi)
%  Lunghezza indicativa: 8-10 pagine.
%  Notebook: 04-14, Simulations*.ipynb.
%========================================================================================

\chapter{Studio delle degenerazioni} \label{chpt:degenerazioni}
\sloppy

\section{Iniezione di un fondo di onde gravitazionali} \label{sec:iniezione_GWB}
% - Tre ampiezze log10_A = -14.5, -15.0, -15.5 con gamma = 13/3, analizzate con il modello
%   a 6 parametri SENZA GWB: il GWB viene assorbito dal red noise?
% - Iniezione con i valori di EPTA DR2new (log10_A = -14, gamma = 2.7) e con gamma = 13/3.
% - Figure candidate: corner_gwb_tre_ampiezze_v2.png, corner_senza_vs_gwb_A145_v2.png, ...,
%   corner_sim_gwb_A140_g27.png, corner_reali_vs_GWB_tre_ampiezze.png.
%   (simulati verde #2A9D8F, con GWB viola #6A4C93)

\section{Degenerazione tra dispersione e scattering} \label{sec:DM_SV}
% - La forma a "L" nei piani (dm_gp, sv_gp): presente già senza GWB o introdotta dal GWB?
% - Esperimento con i valori veri di DM e SV scambiati (swap).

\section{L'indice cromatico dello scattering} \label{sec:alpha}
% - Simulazione con alpha = 4.4, analisi con alpha = 4.0 (modello sbagliato) e 4.4 (corretto).
% - alpha lasciato libero (sv_gp_alpha), con e senza GWB.
% - Dati reali con alpha = 4.4.
% - Figure candidate: corner_sim_sv44_errato_vs_corretto.png, corner_sim_sv44_fixed_vs_alphafree.png,
%   corner_sim_sv44_alphafree_degenerazione.png, corner_gwb_A140_sv44_alphafree_degenerazione.png.

\section{Discussione} \label{sec:discussione}
